\section{Objectifs}
\begin{itemize}
\item Formaliser l’indépendance.
\item Utiliser une partition de l’univers et la formule des probabilités totales.
\item Représenter des épreuves indépendantes par un arbre ou un tableau.
\end{itemize}

\section{Prérequis}
Probabilités conditionnelles, arbres pondérés, évènements.

\section{Activité de découverte}
Comparer $P(A)$ et $P_B(A)$ permet de décider si la connaissance de $B$ modifie la probabilité de $A$.

\section{Vocabulaire}
Indépendance, partition, système complet d’évènements, épreuve de Bernoulli.

\section{Cours}
Deux évènements $A$ et $B$ sont indépendants lorsque
\[
P(A\cap B)=P(A)P(B).
\]

Si $(A_i)$ forme une partition de l’univers, alors
\[
P(B)=\sum_i P(A_i)P_{A_i}(B).
\]

Deux épreuves indépendantes peuvent être représentées par un arbre ou un tableau. Pour une répétition de $n$ épreuves de Bernoulli indépendantes et identiques, le programme demande explicitement le cas $n\le 4$.

\section{Exemples corrigés}
Si $P(A)=0,4$ et $P(B)=0,5$ avec indépendance, alors $P(A\cap B)=0,2$.

\section{Méthode}
Sur un arbre, multiplier le long d’un chemin puis additionner les chemins incompatibles correspondant à l’évènement recherché.

\section{Erreurs fréquentes}
\begin{itemize}
\item Confondre indépendance et incompatibilité.
\item Confondre $P_A(B)$ et $P_B(A)$.
\item Appliquer les probabilités totales sans vérifier la partition.
\end{itemize}

\section{Synthèse}
L’indépendance est une propriété du modèle. La formule des probabilités totales relie les probabilités conditionnelles à une probabilité globale.
